\section{Sensitivity Analysis}

% 系数必须先定义再引用，否则读者无法解释——这是这类章节最常见的漏洞。
Define the relative sensitivity of the index to parameter $p$ as
\begin{equation}
  S_p = \frac{\Delta I / I_0}{\Delta p / p_0},
  \label{eq:sensitivity}
\end{equation}
where $I_0$ is the baseline index at the nominal set and every perturbation is
$\pm 10\,\%$ about it.

\subsection{Sensitivity of Problem 1}
[Rank $|S_p|$; explain any peak by the mechanism, not by the number.]

\subsection{Sensitivity of Problem 2}
[Same procedure on the schedule's parameters.]

\subsection{Sensitivity of Problem 3}
[Observer gains and controller gains; a coefficient that rounds to zero is a finding.]

% 【Section writing prompt】
% 三件事：一个公式（定义敏感性系数）、一个固定摄动幅度（默认 ±10 %，全文统一）、一张排序表。
% 口径不许混：一问答"平均敏感系数"、另一问答"百分比变化"，两套数放一起就没法比较。
% 全为 0 的参数要下结论（"该参数不敏感"），不要罗列五个零。
% 图：三小图并排，纵轴指数、横轴 Parameter Variation (%)、虚线标 Baseline——基线让相对变化可读。
